\documentclass[10pt,a4paper]{amsart}
\usepackage[margin=19mm]{geometry}
\usepackage{amsmath,amssymb,mathtools}
\usepackage[scr=rsfs,cal=euler]{mathalfa}
\usepackage{xfrac}
\usepackage[unicode,hidelinks,bookmarks=false]{hyperref}
\newcommand{\ie}{i.e.\ }
\newcommand{\cf}{cf.\ }
\newcommand{\resp}{respectively\ }
\newcommand{\Subsection}{\subsection}
\providecommand{\nobreakdash}{\nobreak\hbox{-}}
\setlength{\emergencystretch}{2em}
\multlinegap0pt
\usepackage{bm}
\usepackage{bbm}
\usepackage{dsfont}
\usepackage{xspace}
\usepackage{cancel}
\usepackage{mathtools}
\usepackage{amsthm}
\makeatletter
\@ifundefined{Lem}{\newtheorem{Lem}{Lemma}}{}
\@ifundefined{Prop}{\newtheorem{Prop}{Proposition}}{}
\@ifundefined{Thm}{\newtheorem{Thm}{Theorem}}{}
\makeatother
\makeatletter
\expandafter\ifx\csname Mainthm\endcsname\relax
\newenvironment{Mainthm}[1]{%
  \IfBlankTF{#1}
    {\renewcommand{\theMainthminner}{\unskip}}
    {\renewcommand\theMainthminner{#1}}%
  \Mainthminner
}{\endMainthminner}
\fi
\expandafter\ifx\csname Manualthm\endcsname\relax
\newenvironment{Manualthm}[1]{%
  \IfBlankTF{#1}
    {\renewcommand{\theManualthminner}{\unskip}}
    {\renewcommand\theManualthminner{#1}}%
  \Manualthminner
}{\endManualthminner}
\fi
\def\R{\mathbb{R}}
\def\S{\mathbb{S}}
\newcommand{\td}{\mathrm{d}}
\makeatother
\title[Isotropic submanifolds of $T\mathbb{S}^n$ and their focal se]{Isotropic submanifolds of $T\mathbb{S}^n$ and their focal sets}
\author{Nikos Georgiou, Brendan Guilfoyle, Morgan Robson}
\date{Source: arXiv:2601.17451v1 (CC BY 4.0)}
\begin{document}
\maketitle

\noindent\emph{Source and licence.} Isotropic submanifolds of $T\mathbb{S}^n$ and their focal sets. Version arXiv:2601.17451v1 (CC BY 4.0), distributed under the Creative Commons Attribution 4.0 International licence, as recorded in the arXiv licence metadata for this exact version and shown on the abstract page https://arxiv.org/abs/2601.17451v1. Full rights notice: licenses/arxiv-2601.17451-CC-BY-4.0.txt.\par\smallskip

\noindent\textit{Editorial scope.} Counted unit: the Thm whose source label is \texttt{Thm: structure theorem for orthogonal submanifolds} in the article (source0.tex lines 383--389, source label \texttt{Thm: structure theorem for orthogonal submanifolds}), delivered together with the article's own complete original proof of it (source0.tex lines 414--443, one \texttt{proof} environment of 30 lines). No mathematics is written, repaired or re-derived: statement and proof are the article's own text line for line. The proof is detached from the statement in the source: the article states the result at lines 383--389 and prints its proof at lines 414--443, and the proof environment's own header names the statement, so the association is the article's own. Required context retained uncounted: eq:louiville form in terms of round metric (lines 218--220); prop: isotropic implies existance of orthogonal sub manifolds (lines 358--364); Lem: derivative of Phi\_v (lines 391--399). Every definition, display and label of those lines is retained unchanged. Omitted from this package and not redistributed: 1--217, 221--357, 365--382, 390--390, 400--413, 444--1129. Nothing inside a retained excerpt is omitted or altered: every retained body is delivered exactly as the article's lines are written, so no omission receipt of any kind accompanies this package. Citations: the retained excerpts carry no citation command at all, so no bibliography entry of the article is transcribed and no reference is redistributed. Reference adaptation: none was needed and none was made; every reference inside the delivered unit resolves inside it, so no unresolved reference or citation is produced. Printed slips kept literal: no printed word, symbol, hypothesis or number of the counted statement or of its proof was altered. Layout and class: the article's own class is \texttt{amsart} with packages including ulem, soul, amsmath, amsthm, amsfonts, amssymb, amscd, mathrsfs, graphicx, tikz-cd, subfig, comment, fancyvrb, verbatim; the wrapper is the lane's pinned \texttt{amsart} editorial wrapper at 10pt on A4, which restates the theorem-like environments the retained text needs and adds the article's own macro definitions that the retained text needs (\texttt{\textbackslash R}, \texttt{\textbackslash S}, \texttt{\textbackslash td}), with extra wrapper packages (bm, bbm, dsfont, xspace, cancel, mathtools). The delivered numbering is the unit's own. Assets: no retained span contains a figure environment, a graphics-inclusion command, a PDF-page inclusion, a table or a TikZ picture; no image file is copied into this package, the package declares no asset dependency and no image file is redistributed with it. Licence: version arXiv:2601.17451v1 of this article is distributed under the Creative Commons Attribution 4.0 International licence, as recorded in the arXiv licence metadata for this exact version and shown on the abstract page https://arxiv.org/abs/2601.17451v1. A full rights notice is delivered at licenses/arxiv-2601.17451-CC-BY-4.0.txt. Characters: every retained excerpt is pure ASCII, so the delivered engine, XeLaTeX with its default Latin Modern text font, reports no missing character.
\begin{equation}\label{eq:louiville form in terms of round metric}
    \theta_{(p,\beta)}(X)=g_p(\td \pi_{(p,\beta)}(X),\beta),
  \end{equation}

 \begin{Prop} \label{prop: isotropic implies existance of orthogonal sub manifolds}
Let $\Gamma\subseteq T\mathbb{S}^n$ be isotropic and $\dim \Gamma>0$. Suppose that $p\in\Psi(\Gamma\times\mathbb{R})$ is not a focal point. Then $p$ is contained in an orthogonal submanifold of $\Gamma$ of dimension $\dim \Gamma$ which takes the form
\[
\Sigma_0:=\Psi\Big(\Big\{((\xi,\eta),F(\xi,\eta)): (\xi,\eta)\in W\Big\}\Big),
\]
with $W\subset \Gamma$ open and $F:W\to \mathbb{R}$ a smooth function such that $\theta=\td F$ and $F(\xi_0,\eta_0)=r_0$ where $p=\Psi((\xi_0,\eta_0),r_0)$ and $\theta$ is the Liouville 1-form (\ref{eq:louiville form in terms of round metric}).
\end{Prop}

\begin{Lem}\label{Lem: derivative of Phi_v}
    Let $S\subset \R^{n+1}$ be immersed and $v$ a unit section over $S$. For all $X\in T_pS$ the map $\td\Phi_v:TS\to TT\S^n$ satisfies
    \begin{equation}\label{eq: derivative of Phi_v}
    \begin{aligned}
        &\td\pi(\td\Phi_v(X))=\overline\nabla_Xv(p),\\ &\mathcal{K}(\td\Phi_v(X))=X-\langle X,v(p)\rangle v(p)-\langle p, v(p)\rangle \overline\nabla_X v(p),
    \end{aligned}
    \end{equation}
    and $\ker \td\Phi_v|_p \subseteq \mbox{Span}\{v(p)\}$. Furthermore if $v(p)\notin T_pS$, then $\td\Phi_v|_p$ is injective.
\end{Lem}

\begin{Thm}[Structure Theorem for Orthogonal Submanifolds]\label{Thm: structure theorem for orthogonal submanifolds}
    Let $\Gamma\subset T\S^n$ be isotropic and $\mathcal{U}$ an orthogonal submanifold of $\Gamma$ without focal points. For all $p\in \mathcal{U}$ there exists $V_p$, an open neighbourhood of $p$ in $\mathcal{U}$ such that $V_p\subset \Sigma_p$, where
    \[
    \Sigma_p:=\Psi\Big(\Big\{((\xi,\eta),F(\xi,\eta)): (\xi,\eta)\in W\Big\}\Big),
    \]
    with $W\subseteq\Gamma$ open in $\Gamma$ and $F:W\to \mathbb{R}$ a smooth function such that $\theta|_W=\td F$.
\end{Thm}

\begin{proof}[Proof of Theorem \ref{Thm: structure theorem for orthogonal submanifolds}]
    Let $\mathcal{U}$ be an orthogonal submanifold of $\Gamma$ and let $\nu\in\Gamma(N\mathcal{U})$ be unit such that $\Phi_\nu(\mathcal{U})\subseteq \Gamma$. Fix $p\in \mathcal{U}$. Since $\nu(p)\in N_p\mathcal{U}$, the last part of Lemma \ref{Lem: derivative of Phi_v} implies that for some open neighbourhood $V_p\subset \mathcal{U}$ of $p$,  $\Phi_\nu|_{V}$ is a diffeomorphism onto its image.  Define the smooth function
    \begin{align*}
        &r:~\mathcal{U}\to\mathbb{R}, &&q \mapsto \langle q, \nu(q)\rangle.
    \end{align*}
   By the definitions of $\Psi$ and $\Phi_\nu$ one sees that
    $q=\Psi(\Phi_\nu(q),r(q))$ for all $q\in \mathcal{U}$. Hence
    \[
    V_p=\Psi\Big(\Big\{(\Phi_\nu(q),r(q)): q\in V_p\Big\}\Big)=\Psi\Big(\Big\{\big((\xi,\eta),(r\circ\Phi_\nu^{-1})(\xi,\eta)\big): (\xi,\eta)\in \Phi_\nu(V_p)\Big\}\Big).
    \]
    On the other hand, since $\Gamma$ is isotropic and $p$ is not focal, Proposition \ref{prop: isotropic implies existance of orthogonal sub manifolds},
    \[
    \Sigma_p:=\Psi\Big(\Big\{((\xi,\eta),F(\xi,\eta)): (\xi,\eta)\in W\Big\}\Big),
    \]
    is a dimension $\dim \Gamma$ orthogonal submanifold of $\Gamma$ containing $p$, with $\Phi_\nu(p)\in W\subset \Gamma$ open and $F:W\to\mathbb{R}$ a smooth function such that $\td F=\theta$. By shrinking $V_p$ if required we may assume that $\Phi_\nu(V_p)\subset W$. Note further that if $X\in T_q\mathcal{U}$, then
    \begin{align*}
        \td r(X)=X(r)&=\langle X, \nu(q) \rangle + \langle q, \overline\nabla_X \nu(q)\rangle,\\
        &=\langle q, \overline\nabla_X \nu(q)\rangle, &&\text{(Since $X\perp \nu(p)$)}\\
        &=\langle q-\langle q,v(q)\rangle v(q), \overline\nabla_X v(q)\rangle, &&\text{(Since $\nabla_X\nu(p)\perp \nu(p)$)}\\
        &=\langle q-\langle q,v(q)\rangle v(q), \td\pi(\td\Phi_\nu(X))\rangle, &&\text{(By Lemma \ref{Lem: derivative of Phi_v})}\\
        &=\Phi_\nu^\ast\theta(X), &&\text{(From equation (\ref{eq:louiville form in terms of round metric}))}\\
        &=\td(\Phi_\nu^\ast F)(X). &&\text{(By definition of $F$)}
    \end{align*}
    Hence $r=\Phi_\nu^\ast (F|_{\Phi_\nu(V)})$, since we may choose $F$ to satisfy $F(\Phi_\nu(p))=r(p)$. Thus
    \begin{align*}
        V_p=\Psi\Big(\Big\{((\xi,\eta),F(\xi,\eta)):(\xi,\eta)\in \Phi_\nu(V_p)\Big\}\Big)
        \subseteq \Psi\Big(\Big\{((\xi,\eta),F(\xi,\eta)): (\xi,\eta)\in W\Big\}\Big)
        =\Sigma_0.
    \end{align*}
\end{proof}

\end{document}
